\documentclass[11pt,a4paper]{article}
\usepackage[utf8]{inputenc}
\usepackage[T1]{fontenc}
\usepackage[margin=2.5cm]{geometry}
\usepackage{lmodern}
\usepackage{parskip}
\usepackage{microtype}
\usepackage{booktabs}
\usepackage{longtable}
\usepackage{amsmath}
\usepackage{amssymb}
\usepackage{xcolor}
\usepackage{hyperref}
\usepackage{enumitem}
\usepackage{titlesec}
\usepackage{graphicx}
\graphicspath{{docs/figures/}}
\usepackage{multirow}
\usepackage{array}
\usepackage{colortbl}

\hypersetup{
    colorlinks=true,
    linkcolor=blue!60!black,
    urlcolor=blue!60!black,
    citecolor=blue!60!black,
}

\titleformat{\section}{\large\bfseries}{}{0pt}{}[\titlerule]
\titleformat{\subsection}{\normalsize\bfseries}{}{0pt}{}
\titleformat{\subsubsection}{\normalsize\itshape}{}{0pt}{}

% Placeholder command for missing data
\newcommand{\todo}[1]{\colorbox{yellow!40}{\textit{[TODO: #1]}}}
\newcommand{\placeholder}{\multicolumn{1}{c}{\textcolor{gray}{---}}}

\title{\textbf{Wyniki eksperymentów — harmonogramowanie ze wspólnym terminem}\\[4pt]
    \large Porównanie agentów i odniesienie do benchmarku Biskupa i Feldmanna}
\author{Karolina Piotrowska, Jan Rosa}
\date{\today}

\begin{document}
\maketitle
\tableofcontents
\newpage

% ===========================================================================
\section{Konfiguracja eksperymentów}
% ===========================================================================

\subsection{Eksperyment przesiewowy (sweep)}

Eksperyment przesiewowy przeprowadzono w celu porównania czterech agentów
optymalizacyjnych w połączeniu z pięcioma modelami językowymi sterującymi
frameworkiem GEPA.
Parametry stałe dla wszystkich uruchomień:

\begin{center}
\begin{tabular}{@{}ll@{}}
\toprule
\textbf{Parametr} & \textbf{Wartość} \\
\midrule
Zbiór danych              & \texttt{sch1000.txt} (instancje $k = 1\ldots5$) \\
Współczynnik szczelności  & $h = 0{,}4$ \\
Liczba instancji          & 5 (benchmarkowanie) \\
Liczba instancji treningowych & 5 (trening GEPA) \\
Maksymalna liczba kroków  & 5\,000 \\
Budżet wywołań GEPA       & 100 wywołań metryki \\
\bottomrule
\end{tabular}
\end{center}

Porównywane modele LLM (uruchamiane lokalnie przez Ollama):

\begin{center}
\begin{tabular}{@{}ll@{}}
\toprule
\textbf{Model} & \textbf{Klucz LiteLLM} \\
\midrule
Qwen3 4B          & \texttt{ollama/qwen3:4b-instruct-2507-q4\_K\_M} \\
Gemma4 2B         & \texttt{ollama/gemma4:e2b} \\
Phi4-mini         & \texttt{ollama/phi4-mini-reasoning:latest} \\
Granite4 1B       & \texttt{ollama/granite4:1b-h-q8\_0} \\
GPT-OSS 120B      & \texttt{ollama/gpt-oss:120b} \\
\bottomrule
\end{tabular}
\end{center}

\subsection{Eksperyment porównawczy z benchmarkiem}

Eksperyment przeprowadzono w celu odniesienia wyników GEPA-SA do
górnych ograniczeń opublikowanych przez Biskupa i Feldmanna~\cite{biskup2001}.
Parametry:

\begin{center}
\begin{tabular}{@{}ll@{}}
\toprule
\textbf{Parametr} & \textbf{Wartość} \\
\midrule
Zbiór danych              & \texttt{sch1000.txt} (instancje $k = 1\ldots10$) \\
Współczynniki szczelności & $h \in \{0{,}2;\; 0{,}4;\; 0{,}6;\; 0{,}8\}$ \\
Liczba instancji          & 10 \\
Liczba instancji treningowych & 10 \\
Maksymalna liczba kroków  & 5\,000 \\
Budżet wywołań GEPA       & 100 wywołań metryki \\
Agent / model LLM         & GEPA-SA z modelem Granite4~1B \\
\bottomrule
\end{tabular}
\end{center}

Metryka porównawcza to \emph{odchylenie względne} od górnego ograniczenia:
\begin{equation}
  \Delta_{\%}(k, h) = \frac{C^*_{\text{GEPA-SA}}(k, h) - \text{GG}(k, h)}
                           {\text{GG}(k, h)} \times 100\%,
  \label{eq:gap}
\end{equation}
gdzie $C^*_{\text{GEPA-SA}}(k,h)$ to najlepszy koszt znaleziony przez GEPA-SA
na instancji $k$ przy współczynniku $h$, a $\text{GG}(k,h)$ to górne
ograniczenie z publikacji~\cite{biskup2001}.
Wartość ujemna oznacza pobicie opublikowanego górnego ograniczenia.

% ===========================================================================
\section{Eksperyment przesiewowy — porównanie modeli LLM}
% ===========================================================================

\subsection{Zbiorcze zestawienie wyników}

Tabela~\ref{tab:sweep_main} przedstawia średnią procentową poprawę kosztu
względem punktu startowego (inicjalizacja zachłanna) dla każdego agenta
i każdego modelu LLM.
Kolumna $\Delta$ oznacza zysk z zastosowania GEPA w stosunku do bazowej
metaheurystyki (SA lub GA).

\begin{table}[h]
\centering
\caption{Średnia poprawa kosztu [\%] — \texttt{sch1000}, $h=0{,}4$,
         $n=5$ instancji, 5\,000 kroków.}
\label{tab:sweep_main}
\begin{tabular}{@{}lrrrrrrr@{}}
\toprule
\multirow{2}{*}{\textbf{Model LLM}}
  & \multicolumn{3}{c}{\textbf{Wyżarzanie symulowane}}
  & \multicolumn{3}{c}{\textbf{Algorytm genetyczny}} \\
\cmidrule(lr){2-4}\cmidrule(lr){5-7}
  & SA & GEPA-SA & $\Delta$ & Genetic & GEPA-GA & $\Delta$ \\
\midrule
Qwen3 4B         & 63{,}60\% & 63{,}51\% & $-0{,}10$\% &  8{,}78\% & 31{,}49\% & $+22{,}71$\% \\
Granite4 1B      & 62{,}93\% & 63{,}08\% & $+0{,}15$\% &  8{,}98\% & 30{,}34\% & $+21{,}36$\% \\
GPT-OSS 120B     & 62{,}58\% & 62{,}34\% & $-0{,}24$\% &  9{,}36\% & 28{,}57\% & $+19{,}20$\% \\
Phi4-mini        & 62{,}31\% & 62{,}78\% & $+0{,}47$\% & 11{,}91\% & 19{,}76\% &  $+7{,}85$\% \\
Gemma4 2B        & 62{,}71\% & 62{,}09\% & $-0{,}62$\% &  8{,}65\% & 10{,}94\% &  $+2{,}29$\% \\
\midrule
\textbf{Średnia} & \textbf{62{,}83\%} & \textbf{62{,}76\%} & $-0{,}07$\%
                 & \textbf{9{,}54\%}  & \textbf{24{,}22\%} & $+14{,}68$\% \\
\bottomrule
\end{tabular}
\end{table}

\begin{figure}[ht]
\centering
\includegraphics[width=\textwidth]{qwen3_agent_comparison}
\caption{Porównanie agentów dla modelu Qwen3 4B (\texttt{sch1000}, $h=0{,}4$, $n=5$ instancji).
  \emph{Lewy panel}: średnia poprawa kosztu [\%] z odchyleniem standardowym.
  \emph{Prawy panel}: rozkład najlepszego znalezionego kosztu (boxplot; im niżej, tym lepiej).
  SA i GEPA-SA osiągają niemal identyczny poziom, podczas gdy GEPA-GA znacząco
  przewyższa bazowy Genetic, ale z dużą wariancją między instancjami.}
\label{fig:agent_comparison}
\end{figure}

\subsection{Wyżarzanie symulowane (SA i GEPA-SA)}

We wszystkich uruchomieniach SA osiąga średnią poprawę kosztu na poziomie
$62{,}3$–$63{,}6\%$, niezależnie od modelu LLM.
GEPA-SA nie przynosi istotnej poprawy względem bazowego SA
($|\Delta| < 0{,}7\%$ we wszystkich przypadkach), co wskazuje, że
domyślna konfiguracja hiperparametrów wyżarzania jest już bliska optimum
dla danego budżetu kroków.

Godnym uwagi wyjątkiem jest model Gemma4, dla którego GEPA-SA wykonuje
wyraźnie mniej korzystnych ruchów (średnio ok.\ 2\,712 na 5\,000 kroków
wobec $\sim$4\,600 dla pozostałych modeli), ale osiąga zbliżony koszt
końcowy — sugeruje to, że Gemma4 proponuje konfigurację z agresywniejszym
poszukiwaniem, wykonującą rzadsze lecz większe skoki.

\begin{table}[h]
\centering
\caption{Średnia liczba zaakceptowanych ruchów na epizod (SA vs.\ GEPA-SA).}
\label{tab:sa_nimprov}
\begin{tabular}{@{}lrr@{}}
\toprule
\textbf{Model LLM} & \textbf{SA} & \textbf{GEPA-SA} \\
\midrule
Qwen3 4B      & 4\,608 & 4\,611 \\
Granite4 1B   & 4\,634 & 4\,611 \\
GPT-OSS 120B  & 4\,629 & 4\,320 \\
Phi4-mini     & 4\,615 & 4\,621 \\
Gemma4 2B     & 4\,630 & 2\,712 \\
\bottomrule
\end{tabular}
\end{table}

\begin{figure}[ht]
\centering
\includegraphics[width=\textwidth]{qwen3_convergence}
\caption{Krzywe zbieżności kosztu dla wszystkich agentów — Qwen3 4B
  (\texttt{sch1000}, $h=0{,}4$, $n=5$ instancji, 5\,000 kroków).
  Linia ciągła to średnia po instancjach, pasmo to $\pm 1\sigma$,
  linia przerywana to najlepszy wynik spośród instancji.
  SA i GEPA-SA zbiegają gwałtownie w pierwszych $\sim$500 krokach
  i osiągają stabilny poziom kosztu $\approx 7{,}3 \times 10^6$.
  GEPA-GA zbiega wolniej i ze znacznie wyższą wariancją.
  Bazowy Genetic praktycznie nie poprawia kosztu po pierwszych 1000 krokach.}
\label{fig:convergence}
\end{figure}

\subsection{Algorytm genetyczny (Genetic i GEPA-GA)}

Bazowy algorytm genetyczny z domyślnymi hiperparametrami osiąga jedynie
$8{,}7$–$11{,}9\%$ poprawy, co czyni go znacznie słabszym od SA.
GEPA-GA poprawia ten wynik istotnie, ale jakość zależy silnie od modelu LLM.

\begin{table}[h]
\centering
\caption{Szczegóły GEPA-GA: zakres i zmienność poprawy kosztu
         (\texttt{sch1000}, $h=0{,}4$, $n=5$ instancji).}
\label{tab:gepaga_detail}
\begin{tabular}{@{}lrrrr@{}}
\toprule
\textbf{Model LLM} & \textbf{min} & \textbf{śr.} & \textbf{max}
                   & \textbf{śr.\ $n_{\text{impv}}$} \\
\midrule
Qwen3 4B      &  3{,}75\% & 31{,}49\% & 42{,}75\% &   943 \\
Granite4 1B   & 14{,}56\% & 30{,}34\% & 46{,}16\% &   811 \\
GPT-OSS 120B  &  2{,}37\% & 28{,}57\% & 49{,}25\% &   724 \\
Phi4-mini     & 16{,}57\% & 19{,}76\% & 22{,}95\% &   251 \\
Gemma4 2B     &  5{,}46\% & 10{,}94\% & 14{,}79\% &   140 \\
\bottomrule
\end{tabular}
\end{table}

\begin{figure}[ht]
\centering
\includegraphics[width=\textwidth]{qwen3_gepa_ga_history}
\caption{Historia optymalizacji GEPA-GA przy modelu Qwen3 4B (32 wywołania metryki).
  \emph{Górny lewy}: wynik złożony (score) każdego wywołania i kumulatywne optimum.
  \emph{Górny prawy}: rozrzut poprawy kosztu po instancjach dla kolejnych wywołań.
  \emph{Dolny lewy}: ewolucja hiperparametrów GA (wartości unormowane do $[0,1]$).
  Już w pierwszej iteracji GEPA skacze z konfiguracji domyślnej ($\sim$3\%) do
  $\sim$56\% poprawy kosztu — znajdując \texttt{population\_size=50}, \texttt{generations=10},
  co odpowiada agresywnemu przeszukiwaniu populacyjnemu.}
\label{fig:gepa_ga_qwen3}
\end{figure}

\begin{figure}[ht]
\centering
\includegraphics[width=\textwidth]{gemma4_gepa_ga_history}
\caption{Historia optymalizacji GEPA-GA przy modelu Gemma4 2B (32 wywołania metryki) —
  przypadek kontrastowy wobec Rys.~\ref{fig:gepa_ga_qwen3}.
  Skala osi Y jest znacznie węższa (0--20\% vs.\ 0--60\%):
  Gemma4 nie jest w stanie wyprowadzić GA poza $\sim$14\% poprawy.
  Hiperparametry oscylują blisko wartości domyślnych przez cały przebieg optymalizacji,
  co wskazuje na niezdolność modelu do zaproponowania jakościowo nowych konfiguracji.}
\label{fig:gepa_ga_gemma4}
\end{figure}

Obserwacje:

\begin{itemize}
  \item \textbf{Qwen3, Granite4, GPT-OSS} — modele te proponują konfiguracje
        dające wysoką średnią poprawę ($\sim$29–31\%), ale z dużą wariancją
        między instancjami (nawet $2$–$5\%$ dla niektórych).
        Sugeruje to wrażliwość znalezionych konfiguracji GA na konkretną
        instancję problemu.

  \item \textbf{Phi4-mini} — najniższa wariancja ($16$–$23\%$) kosztem
        niższego sufitu.  Model proponuje zachowawcze, stabilne konfiguracje.

  \item \textbf{Gemma4} — efekt GEPA jest znikomy ($+2{,}3$ pp.\ ponad
        bazowy GA).  Model nie jest w stanie wygenerować konfiguracji
        wyraźnie lepszych od wartości domyślnych.
\end{itemize}

\subsection{Wnioski z eksperymentu przesiewowego}

\begin{enumerate}[label=\arabic*.]
  \item SA zdecydowanie dominuje nad GA przy 5\,000 krokach —
        różnica wynosi ponad 30 punktów procentowych.
  \item GEPA nie przynosi zysku dla SA (domyślna konfiguracja jest
        wystarczająco dobra); dla GA zysk jest znaczący i zależny od modelu.
  \item Qwen3 4B i Granite4 1B oferują najlepszy stosunek jakości do
        zasobów (oba modele $\leq 2$ B parametrów, lokalne, bez opóźnień API).
  \item GPT-OSS 120B pomimo rozmiaru nie przewyższa Qwen3 4B, a wykazuje
        wyższą wariancję wyników GEPA-GA.
\end{enumerate}

% ===========================================================================
\section{Porównanie z benchmarkiem Biskupa i Feldmanna}
% ===========================================================================

\subsection{Górne ograniczenia z publikacji}

Tabela~\ref{tab:bishop_ub} zawiera górne ograniczenia wartości funkcji celu
dla 10 instancji o $n = 1000$ zadaniach, wyznaczone przez Biskupa i Feldmanna
heurystyką~II~\cite{biskup2001} (Tabela~9 w publikacji).
Wartości te są górną granicą optymalnego rozwiązania — algorytm, który
osiągnie koszt niższy od podanego, poprawia dotychczas znany wynik.

\begin{table}[h]
\centering
\caption{Górne ograniczenia dla $n = 1000$ (Tabela~9 z~\cite{biskup2001}).
         Wszystkie wartości uzyskane heurystyką~II.}
\label{tab:bishop_ub}
\begin{tabular}{@{}rrrrr@{}}
\toprule
$k$ & $h = 0{,}2$ & $h = 0{,}4$ & $h = 0{,}6$ & $h = 0{,}8$ \\
\midrule
 1 & 15\,190\,371 &  8\,570\,154 & 6\,411\,581 & 6\,411\,581 \\
 2 & 13\,356\,727 &  7\,592\,040 & 6\,112\,598 & 6\,112\,598 \\
 3 & 12\,919\,259 &  7\,313\,736 & 5\,985\,538 & 5\,985\,538 \\
 4 & 12\,705\,290 &  7\,300\,217 & 6\,096\,729 & 6\,096\,729 \\
 5 & 13\,276\,868 &  7\,738\,367 & 6\,348\,242 & 6\,348\,242 \\
 6 & 12\,236\,080 &  7\,144\,491 & 6\,082\,142 & 6\,082\,142 \\
 7 & 14\,160\,773 &  8\,426\,024 & 6\,575\,879 & 6\,575\,879 \\
 8 & 13\,314\,723 &  7\,508\,507 & 6\,069\,658 & 6\,069\,658 \\
 9 & 12\,433\,821 &  7\,299\,271 & 6\,188\,416 & 6\,188\,416 \\
10 & 13\,395\,234 &  7\,617\,658 & 6\,147\,295 & 6\,147\,295 \\
\midrule
\textbf{Śr.}
   & 13\,298\,945 &  7\,651\,047 & 6\,201\,808 & 6\,201\,808 \\
\bottomrule
\end{tabular}
\end{table}

Uwaga: dla $h = 0{,}6$ i $h = 0{,}8$ wartości górnych ograniczeń są
identyczne dla każdej instancji.
Jest to wynik szczelności terminu — przy wystarczająco luźnym terminie
optymalny harmonogram staje się niezależny od dokładnej wartości $h$.

\subsection{Wyniki GEPA-SA}

Tabele~\ref{tab:h02}--\ref{tab:h08} zawierają najlepszy koszt znaleziony przez
GEPA-SA dla każdej z 10 instancji oraz odchylenie względne $\Delta_{\%}$
(wzór~\ref{eq:gap}) od górnego ograniczenia Biskupa i Feldmanna.
Wartość ujemna oznacza pobicie opublikowanego ograniczenia.
Użyto modelu Granite4~1B — najlepszego modelu w eksperymencie przesiewowym
pod względem poprawy GEPA-GA; dla samego GEPA-SA różnice między modelami
są poniżej $0{,}5$ pp., więc wybór modelu nie ma tu istotnego wpływu.

% ---- h = 0.2 ---------------------------------------------------------------
\begin{table}[ht]
\centering
\caption{Wyniki dla $h = 0{,}2$ — \texttt{sch1000}, 5\,000 kroków, Granite4~1B.}
\label{tab:h02}
\begin{tabular}{@{}rrrr@{}}
\toprule
$k$ & GG (Biskup) & $C^*_{\text{GEPA-SA}}$ & $\Delta_{\text{GEPA-SA}}$ \\
\midrule
 1 & 15\,190\,371 & 14\,077\,332 & $-7{,}33\%$ \\
 2 & 13\,356\,727 & 12\,318\,195 & $-7{,}78\%$ \\
 3 & 12\,919\,259 & 11\,991\,068 & $-7{,}18\%$ \\
 4 & 12\,705\,290 & 11\,819\,056 & $-6{,}98\%$ \\
 5 & 13\,276\,868 & 12\,475\,975 & $-6{,}03\%$ \\
 6 & 12\,236\,080 & 11\,669\,284 & $-4{,}63\%$ \\
 7 & 14\,160\,773 & 13\,299\,127 & $-6{,}08\%$ \\
 8 & 13\,314\,723 & 12\,295\,773 & $-7{,}65\%$ \\
 9 & 12\,433\,821 & 11\,782\,948 & $-5{,}23\%$ \\
10 & 13\,395\,234 & 12\,450\,882 & $-7{,}05\%$ \\
\midrule
\textbf{Śr.} & 13\,298\,915 & 12\,417\,964 & $\mathbf{-6{,}62\%}$ \\
\bottomrule
\end{tabular}
\end{table}

% ---- h = 0.4 ---------------------------------------------------------------
\begin{table}[ht]
\centering
\caption{Wyniki dla $h = 0{,}4$ — \texttt{sch1000}, 5\,000 kroków, Granite4~1B.}
\label{tab:h04}
\begin{tabular}{@{}rrrr@{}}
\toprule
$k$ & GG (Biskup) & $C^*_{\text{GEPA-SA}}$ & $\Delta_{\text{GEPA-SA}}$ \\
\midrule
 1 &  8\,570\,154 &  8\,128\,044 & $-5{,}16\%$ \\
 2 &  7\,592\,040 &  7\,293\,467 & $-3{,}93\%$ \\
 3 &  7\,313\,736 &  7\,009\,299 & $-4{,}16\%$ \\
 4 &  7\,300\,217 &  7\,055\,303 & $-3{,}35\%$ \\
 5 &  7\,738\,367 &  7\,385\,882 & $-4{,}56\%$ \\
 6 &  7\,144\,491 &  6\,951\,632 & $-2{,}70\%$ \\
 7 &  8\,426\,024 &  7\,886\,057 & $-6{,}41\%$ \\
 8 &  7\,508\,507 &  7\,256\,262 & $-3{,}36\%$ \\
 9 &  7\,299\,271 &  7\,086\,209 & $-2{,}92\%$ \\
10 &  7\,617\,658 &  7\,309\,020 & $-4{,}05\%$ \\
\midrule
\textbf{Śr.} & 7\,651\,047 & 7\,336\,118 & $\mathbf{-4{,}12\%}$ \\
\bottomrule
\end{tabular}
\end{table}

% ---- h = 0.6 ---------------------------------------------------------------
\begin{table}[ht]
\centering
\caption{Wyniki dla $h = 0{,}6$ — \texttt{sch1000}, 5\,000 kroków, Granite4~1B.}
\label{tab:h06}
\begin{tabular}{@{}rrrr@{}}
\toprule
$k$ & GG (Biskup) & $C^*_{\text{GEPA-SA}}$ & $\Delta_{\text{GEPA-SA}}$ \\
\midrule
 1 & 6\,411\,581 & 6\,445\,488 & $+0{,}53\%$ \\
 2 & 6\,112\,598 & 6\,187\,801 & $+1{,}23\%$ \\
 3 & 5\,985\,538 & 6\,099\,638 & $+1{,}91\%$ \\
 4 & 6\,096\,729 & 6\,224\,106 & $+2{,}09\%$ \\
 5 & 6\,348\,242 & 6\,479\,263 & $+2{,}06\%$ \\
 6 & 6\,082\,142 & 6\,251\,208 & $+2{,}78\%$ \\
 7 & 6\,575\,879 & 6\,639\,210 & $+0{,}96\%$ \\
 8 & 6\,069\,658 & 6\,306\,187 & $+3{,}90\%$ \\
 9 & 6\,188\,416 & 6\,182\,658 & $-0{,}09\%$ \\
10 & 6\,147\,295 & 6\,250\,481 & $+1{,}68\%$ \\
\midrule
\textbf{Śr.} & 6\,201\,808 & 6\,306\,604 & $\mathbf{+1{,}69\%}$ \\
\bottomrule
\end{tabular}
\end{table}

% ---- h = 0.8 ---------------------------------------------------------------
\begin{table}[ht]
\centering
\caption{Wyniki dla $h = 0{,}8$ — \texttt{sch1000}, 5\,000 kroków, Granite4~1B.}
\label{tab:h08}
\begin{tabular}{@{}rrrr@{}}
\toprule
$k$ & GG (Biskup) & $C^*_{\text{GEPA-SA}}$ & $\Delta_{\text{GEPA-SA}}$ \\
\midrule
 1 & 6\,411\,581 & 8\,609\,536 & $+34{,}28\%$ \\
 2 & 6\,112\,598 & 8\,507\,317 & $+39{,}18\%$ \\
 3 & 5\,985\,538 & 8\,729\,646 & $+45{,}85\%$ \\
 4 & 6\,096\,729 & 8\,785\,082 & $+44{,}10\%$ \\
 5 & 6\,348\,242 & 9\,199\,654 & $+44{,}92\%$ \\
 6 & 6\,082\,142 & 9\,060\,791 & $+48{,}97\%$ \\
 7 & 6\,575\,879 & 9\,145\,941 & $+39{,}08\%$ \\
 8 & 6\,069\,658 & 8\,652\,371 & $+42{,}55\%$ \\
 9 & 6\,188\,416 & 8\,771\,697 & $+41{,}74\%$ \\
10 & 6\,147\,295 & 8\,804\,172 & $+43{,}22\%$ \\
\midrule
\textbf{Śr.} & 6\,201\,808 & 8\,826\,621 & $\mathbf{+42{,}32\%}$ \\
\bottomrule
\end{tabular}
\end{table}

\subsection{Zestawienie zbiorcze}

Tabela~\ref{tab:summary} zbiera średnie wyniki dla wszystkich czterech
wartości $h$.

\begin{table}[ht]
\centering
\caption{Zestawienie zbiorcze — średnie wyniki po 10 instancjach
         (\texttt{sch1000}, 5\,000 kroków, Granite4~1B).}
\label{tab:summary}
\begin{tabular}{@{}lrrr@{}}
\toprule
$h$ & GG śr.\ (Biskup) & $C^*$ śr.\ GEPA-SA & $\Delta$ GEPA-SA \\
\midrule
0{,}2 & 13\,298\,915 & 12\,417\,964 & $-6{,}62\%$ \\
0{,}4 &  7\,651\,047 &  7\,336\,118 & $-4{,}12\%$ \\
0{,}6 &  6\,201\,808 &  6\,306\,604 & $+1{,}69\%$ \\
0{,}8 &  6\,201\,808 &  8\,826\,621 & $+42{,}32\%$ \\
\bottomrule
\end{tabular}
\end{table}

% ===========================================================================
\section{Wnioski}
% ===========================================================================

\subsection{Wnioski z eksperymentu przesiewowego}

Wyżarzanie symulowane (SA) okazało się dominującą metodą dla instancji
\texttt{sch1000} przy 5\,000 krokach — niezależnie od modelu LLM i bez
istotnej pomocy ze strony GEPA.
Główna wartość frameworku GEPA ujawnia się przy algorytmie genetycznym,
gdzie automatyczny dobór hiperparametrów pozwolił potroić skuteczność
najlepszych modeli (Qwen3, Granite4) w porównaniu z konfiguracją domyślną.

Spośród badanych modeli LLM Qwen3 4B i Granite4 1B wyróżniają się najlepszym
stosunkiem uzyskanej jakości do kosztu obliczeniowego (oba modele lokalne,
$\leq 2$B parametrów).

\subsection{Wnioski z porównania z benchmarkiem}

Wyniki zestawione w Tabelach~\ref{tab:h02}--\ref{tab:h08} prowadzą do
następujących wniosków:

\begin{enumerate}[label=\arabic*.]
  \item \textbf{Napięty termin ($h = 0{,}2$ i $h = 0{,}4$) — GEPA-SA
        pokonuje górne ograniczenia Biskupa.}
        Przy $h = 0{,}2$ GEPA-SA osiąga średnio $-6{,}62\%$ poniżej opublikowanego
        górnego ograniczenia (zakres $-4{,}63\%$ do $-7{,}78\%$).
        Przy $h = 0{,}4$ analogicznie $-4{,}12\%$ (zakres $-2{,}70\%$ do $-6{,}41\%$).
        Oznacza to, że wyżarzanie symulowane przy 5\,000 krokach jest silniejszą
        heurystyką niż obie metody zaproponowane przez Biskupa i Feldmanna
        dla napiętych terminów.

  \item \textbf{Umiarkowanie luźny termin ($h = 0{,}6$) — GEPA-SA nieznacznie
        powyżej górnego ograniczenia.}
        Średnia luka wynosi $+1{,}69\%$ — GEPA-SA jest o ok.\ 2\% gorszy od
        heurystyki Biskupa.  Różnica jest mała, ale jednolita dla wszystkich
        10 instancji (jedyny wyjątek: $k = 9$, gdzie GEPA-SA uzyskuje $-0{,}09\%$).
        Wynik sugeruje, że SA gorzej radzi sobie ze strukturą V-kształtną
        optymalnego harmonogramu przy luźnym terminie.

  \item \textbf{Bardzo luźny termin ($h = 0{,}8$) — SA znacząco poniżej
        górnego ograniczenia.}
        Luka wynosi średnio $+42{,}32\%$ — SA daje koszty prawie dwukrotnie
        wyższe od wartości opublikowanych.
        Zwraca uwagę, że górne ograniczenia Biskupa dla $h = 0{,}8$ są
        identyczne jak dla $h = 0{,}6$ (własność nierestrykacyjności terminu),
        podczas gdy nasze wyniki dla $h = 0{,}8$ są drastycznie gorsze niż
        dla $h = 0{,}6$.
        Wskazuje to, że SA z bieżącą inicjalizacją zachłanną i harmonogramem
        temperatury nie jest przystosowany do optymalizacji przy bardzo luźnych
        terminach — dominują kary za wczesność, a przestrzeń poszukiwań
        wymaga innej strategii eksploracji.

  \item \textbf{Dla $h = 0{,}8$ problem leży w algorytmie, nie w hiperparametrach.}
        Wyniki z eksperymentu przesiewowego pokazują, że GEPA-SA i bazowy SA
        osiągają niemal identyczne koszty (różnica $< 0{,}1$ pp.)\ niezależnie
        od wartości $h$.
        GEPA nie jest w stanie skompensować strukturalnego ograniczenia SA
        przy bardzo luźnych terminach — konieczna jest zmiana samego algorytmu
        lub strategii inicjalizacji.
\end{enumerate}

\begin{thebibliography}{1}

\bibitem{biskup2001}
D.~Biskup, M.~Feldmann,
\textit{Benchmarks for scheduling on a single machine against restrictive and
unrestrictive common due dates},
Computers \& Operations Research, vol.~28, pp.~787--801, 2001.

\end{thebibliography}

\end{document}
